% Provenance (never printed): sources/book2/BRIEFS.md, g12-09. The consumption function; planned expenditure and
% equilibrium output; the multiplier as a geometric series, with taxes and imports; the toy model and what shrinks the
% multiplier; measured multipliers (debate). Real numbers: figdata/grade-12/09-keynesian-cross/ (Ramey 2019 Table 1 as a
% cited table; Brenland's propensities fitted on the invented accounts). Ledger: Ramey 2019 (abstract, text, Tables 1 and 3),
% Blanchard-Leigh 2013 abstract, and Book 3's rows on Ramey-Zubairy 2018, Auerbach-Gorodnichenko 2012 and
% Nakamura-Steinsson 2014.
\chapter{The Keynesian Cross and the Multiplier}\label{ch:g12:keynesian-cross}

When a government spends one more dollar, by how much does output rise? Since the financial crisis of 2008, economists have estimated the answer in dozens of studies,
with different methods, countries and periods. A survey of them found that most estimates of the average government-spending multiplier lie between 0.6 and 1: one more
dollar of \omterm{def:g10:national-accounts:investment}{government purchases} raises output by less than a dollar, or by about one.
% ledger: g12:ramey-abstract

This chapter builds the simplest model of the question, in which the answer can be much larger, then shows each assumption that makes it large, and what happens to the
answer when each is dropped.

\begin{openingchart}{Estimates of the cumulative government-spending multiplier, from time-series studies (blue) and estimated models of the whole economy (red); a dot is a
single estimate, a bar the range a study reports across samples or horizons.}{Ramey (2019), Table 1, and the studies it surveys}
\begin{tikzpicture}
\begin{axis}[omscatter, width=0.62\linewidth, height=7.4cm, xmin=0, xmax=2.2, ymin=0.3, ymax=14.7, xlabel={multiplier},
  ytick={1,...,14}, yticklabels from table={figdata/grade-12/09-keynesian-cross/out/multipliers.csv}{label}, y dir=reverse,
  y tick label style={font=\scriptsize}, xtick={0,0.5,...,2}, ymajorgrids=false, xmajorgrids=true]
\addplot[gray, dashed, sharp plot, no marks, samples=2, domain=0.3:14.7] ({1}, {x});
\addplot[only marks, mark=*, mark size=1.6pt, color=omDef, error bars/.cd, x dir=both, x explicit, error bar style={line width=1.2pt}]
  table[x=mid, y expr=15-\thisrow{pos}, x error minus=em, x error plus=ep, restrict expr to domain={\thisrow{pos}}{6:14}]
  {figdata/grade-12/09-keynesian-cross/out/multipliers.csv};
\addplot[only marks, mark=*, mark size=1.6pt, color=omThm, error bars/.cd, x dir=both, x explicit, error bar style={line width=1.2pt}]
  table[x=mid, y expr=15-\thisrow{pos}, x error minus=em, x error plus=ep, restrict expr to domain={\thisrow{pos}}{1:5}]
  {figdata/grade-12/09-keynesian-cross/out/multipliers.csv};
\end{axis}
\end{tikzpicture}
\end{openingchart}

\begin{recall}
From \cref{ch:g10:national-accounts}: output equals spending, $Y = C + I + G + X - M$. From \cref{ch:g12:saving-and-investment}: \omterm{def:g12:saving-and-investment:crowding-out}{crowding out}. From
\cref{ch:g12:aggregate-demand-and-supply}: \omterm{def:g12:aggregate-demand-and-supply:ad}{aggregate demand}, \omterm{def:g12:aggregate-demand-and-supply:sticky-price}{sticky prices} and \omterm{def:g12:aggregate-demand-and-supply:potential-output}{potential output}. From mathematics: the sum of a geometric series.
\end{recall}

\section{Planned spending and equilibrium output}\label{sec:g12:keynesian-cross:equilibrium}

\begin{definition}[Consumption function]\label{def:g12:keynesian-cross:consumption-function}
A \emph{consumption function}\index{consumption function} gives households' consumption as a function of their disposable income: here $C = C_0 + c\,(Y - T)$, where $C_0$
is consumption that does not depend on income.
\end{definition}

\begin{definition}[Marginal propensity to consume]\label{def:g12:keynesian-cross:mpc}
The \emph{marginal propensity to consume}\index{marginal propensity to consume} $c$ is the share of an extra unit of disposable income that is spent on consumption.
\end{definition}

\begin{definition}[Planned expenditure and autonomous spending]\label{def:g12:keynesian-cross:planned-expenditure}
\emph{Planned expenditure}\index{planned expenditure} is the total that households, firms, the government and foreigners plan to spend at a given level of output.
\emph{Autonomous spending}\index{autonomous spending}\phantomsection\label{def:g12:keynesian-cross:autonomous} is the part of it that does not depend on output.
\end{definition}

In a closed economy without taxes, \omterm{def:g12:keynesian-cross:planned-expenditure}{planned expenditure} is $PE = C_0 + cY + I + G$. If firms produce more than is planned to be bought, unsold goods pile up and they cut
output; if they produce less, stocks run down and they raise it. Output settles where $Y = PE$. With $C_0 = 100$, $c = 0.8$, $I = 50$ and $G = 50$, \omterm{def:g12:keynesian-cross:planned-expenditure}{autonomous spending} is
200 and $Y = 200 + 0.8Y$ gives $Y = 1000$. If government spending rises by 100, $Y = 300 + 0.8Y$ gives $Y = 1500$: output rises by 500, five times the rise in spending.

\begin{omfigure}
\begin{tikzpicture}
\begin{axis}[width=0.62\linewidth, height=6.0cm, xmin=0, xmax=2000, ymin=0, ymax=2000, xlabel={output $Y$}, ylabel={planned expenditure},
  xtick={400,1000,1500}, ytick={400,1000,1500}, grid=major, grid style={gray!20}, scaled ticks=false]
\addplot[gray, thick, domain=0:2000, samples=2] {x};
\addplot[omDef, thick, domain=0:2000, samples=2] {200 + 0.8*x};
\addplot[omDef, thick, dashed, domain=0:2000, samples=2] {300 + 0.8*x};
\addplot[omThm, thick, dashed, domain=0:2000, samples=2] {200 + 0.5*x};
\node[font=\scriptsize, anchor=north west] at (axis cs:1450,1800) {$PE = Y$};
\node[font=\scriptsize, omDef, anchor=south east] at (axis cs:600,780) {$PE'$};
\node[font=\scriptsize, omDef, anchor=north west] at (axis cs:600,680) {$PE$};
\node[font=\scriptsize, omThm, anchor=north west] at (axis cs:1300,840) {with leakages};
\fill (axis cs:1000,1000) circle (1.5pt) (axis cs:1500,1500) circle (1.5pt) (axis cs:400,400) circle (1.5pt);
\end{axis}
\end{tikzpicture}

{\small The cross. \omterm{def:g12:keynesian-cross:planned-expenditure}{Planned expenditure} $PE = 200 + 0.8Y$ meets the line $PE = Y$ at 1000; a rise of 100 in spending moves it to 1500. With taxes and imports (the red line,
slope 0.5), the same \omterm{def:g12:keynesian-cross:planned-expenditure}{autonomous spending} gives output 400, and each extra unit of spending raises output by 2, not 5.}
\end{omfigure}

\section{The multiplier}\label{sec:g12:keynesian-cross:multiplier}

\begin{definition}[Spending multiplier]\label{def:g12:keynesian-cross:multiplier}
The \emph{spending multiplier}\index{spending multiplier} is the change in output caused by a change in \omterm{def:g12:keynesian-cross:planned-expenditure}{autonomous spending}, divided by that change.
\end{definition}

\begin{proposition}[The multiplier in the cross]\label{prop:g12:keynesian-cross:multiplier}
In the model without taxes or imports, with $0 < c < 1$, a rise $\Delta A$ in \omterm{def:g12:keynesian-cross:planned-expenditure}{autonomous spending} raises output by $\Delta A / (1 - c)$.
\end{proposition}
\begin{proof}
The rise $\Delta A$ is someone's income; they spend $c\,\Delta A$ of it, which is someone else's income, who spends $c^2 \Delta A$, and so on. After $n$ rounds output has
risen by $\Delta A (1 + c + \dots + c^{n-1}) = \Delta A (1 - c^n)/(1 - c)$. Since $0 < c < 1$, $c^n \to 0$, and the total tends to $\Delta A/(1 - c)$. The same follows in one
line from $Y = A + cY$: $Y = A/(1 - c)$, so $\Delta Y = \Delta A/(1 - c)$.
\end{proof}

\begin{omfigure}
\begin{tikzpicture}
\begin{axis}[ybar, /pgf/bar width=10pt, width=0.62\linewidth, height=4.6cm, xmin=0.4, xmax=10.6, ymin=0, ymax=110, xlabel={round of spending},
  ylabel={new spending}, xtick={1,...,10}, grid=major, grid style={gray!20}]
\addplot[fill=omDef!60, draw=omDef, domain=1:10, samples=10] {100*0.8^(x-1)};
\end{axis}
\end{tikzpicture}

{\small A rise of 100 in spending with $c = 0.8$: each round spends 80\,\% of the one before. After three rounds output has risen by 244, after ten by 446; the total
tends to 500.}
\end{omfigure}

\begin{definition}[Marginal propensity to import]\label{def:g12:keynesian-cross:mpm}
The \emph{marginal propensity to import}\index{marginal propensity to import} $m$ is the share of an extra unit of output that is spent on imports.
\end{definition}

\begin{proposition}[The multiplier with taxes and imports]\label{prop:g12:keynesian-cross:leakages}
If taxes are a share $t$ of output and imports a share $m$, the multiplier is $1/(1 - c(1 - t) + m)$, smaller than $1/(1 - c)$ when $t > 0$ or $m > 0$.
\end{proposition}
\begin{proof}
Consumption is $C_0 + c(1 - t)Y$ and imports are $mY$, so $Y = A + c(1 - t)Y - mY$ and $Y = A/(1 - c(1 - t) + m)$. The denominator rises with $t$ and with $m$.
\end{proof}

Taxes and imports are leaks: each round of spending loses part of its income to the government and part to foreign producers, so the rounds shrink faster.

\section{A toy model}\label{sec:g12:keynesian-cross:toy}

\begin{model}{The Keynesian cross}
\index{Keynesian cross}%
\assumes Prices are fixed; the interest rate is fixed; households spend a fixed share of their current income and do not look ahead; firms have idle capacity and produce
  whatever is demanded.
\predicts A rise in \omterm{def:g12:keynesian-cross:planned-expenditure}{autonomous spending} raises output by the multiplier $1/(1 - c(1 - t) + m)$.
\tested Measured
  multipliers are mostly between 0.6 and 1, some higher in deep slumps or when interest rates are stuck at zero, lower with a floating exchange rate and open trade. A large,
  fixed multiplier, the same in all conditions, is refuted by the record.
\breaks At full capacity, where more demand raises prices rather than output; when the central bank raises interest rates in response; and when households look ahead and
  expect higher taxes later.
\end{model}
% ledger: g12:ramey-abstract, g12:ramey-exchange

Each assumption removes a leak. If prices can rise, part of the extra demand goes into prices, as in \cref{ch:g12:aggregate-demand-and-supply}. If the central bank raises
its rate, investment falls, the \omterm{def:g12:saving-and-investment:crowding-out}{crowding out} of \cref{ch:g12:saving-and-investment}. If households know that extra spending today means taxes tomorrow, they spend less of
a temporary rise in income. And at full capacity, output cannot rise at all.

\begin{debate}[How large is the government-spending multiplier?]
% ledger: b1:ramey2018, b1:auerbach2012, b1:nakamura2014, g12:bl2013, g12:ramey-br, g12:ramey-exchange, g12:ramey-chodorow
\emph{Below one, in good times and bad.} Ramey and Zubairy (2018), with quarterly American data reaching back more than a century, find multipliers below one whether or not
the economy has slack; at the zero lower bound the results are mixed, a few specifications as high as 1.5. Their shocks come mostly from military spending in wars, which may
differ from other spending.

\emph{Larger in \omterm{def:g12:aggregate-demand-and-supply:recession}{recessions}.} Auerbach and Gorodnichenko (2012), with models whose parameters switch between expansions and \omterm{def:g12:aggregate-demand-and-supply:recession}{recessions}, find \omterm{def:g12:keynesian-cross:multiplier}{spending multipliers} much larger in
\omterm{def:g12:aggregate-demand-and-supply:recession}{recessions}. The finding depends on how the states are defined and on the sample.

\emph{Regional evidence.} Nakamura and Steinsson (2014), comparing American regions during military build-ups, find a relative multiplier of about 1.5; cross-state studies of
the 2009 stimulus give 1.7 to 2. But a region shares its country's interest rate and taxes, so a relative multiplier is not the national one.

\emph{Forecast errors.} Blanchard and Leigh (2013) found that, early in the crisis, countries that planned larger budget cuts grew less than forecast: multipliers were
substantially higher than forecasters had assumed. The method gives no precise value.

\emph{Annual defence spending.} Regressions on annual American defence spending give 0.6 to 0.7; the variation comes mostly from a few wars.

\emph{Exchange rates.} Across 44 countries, multipliers reach about 1.4 in the long run under fixed exchange rates, and lie between 0.1 and $-0.7$ under flexible ones
(Ilzetzki, Mendoza and V\'egh 2010). Countries choose their regime, so the split between fixed and flexible is not random.
\end{debate}

\section{Exercises}

\begin{exercise}[$\star$]\label{exo:g12:keynesian-cross:1}
A closed economy without taxes has $C = 100 + 0.8Y$, $I = 100$ and $G = 100$. Find equilibrium output.
\end{exercise}

\begin{exercise}[$\star$]\label{exo:g12:keynesian-cross:2}
$C = 50 + 0.75Y$, $I = 100$, $G = 50$, no taxes or trade. Find equilibrium output.
\end{exercise}

\begin{exercise}[$\star$]\label{exo:g12:keynesian-cross:3}
Find the multiplier $1/(1 - c)$ for $c = 0.6$ and for $c = 0.9$.
\end{exercise}

\begin{exercise}[$\star$]\label{exo:g12:keynesian-cross:4}
From the opening chart, how many of the fourteen estimates reach above 1?
\end{exercise}

\begin{exercise}[$\star$]\label{exo:g12:keynesian-cross:5}
From the opening chart, which estimates are the lowest, and from what kind of event do they come?
\end{exercise}

\begin{exercise}[$\star\star$]\label{exo:g12:keynesian-cross:6}
Find the multiplier for $c = 0.8$, $t = 0.25$ and $m = 0.1$.
\end{exercise}

\begin{exercise}[$\star\star$]\label{exo:g12:keynesian-cross:7}
Find the multiplier for $c = 0.9$, $t = 0.3$ and $m = 0.3$, and compare it with $1/(1 - c)$.
\end{exercise}

\begin{exercise}[$\star\star$]\label{exo:g12:keynesian-cross:8}
With the multiplier of \cref{exo:g12:keynesian-cross:6}, government spending rises by 50. By how much does output rise in the model? By how much do tax receipts rise, and
by how much the deficit?
\end{exercise}

\begin{exercise}[$\star\star$]\label{exo:g12:keynesian-cross:9}
With $c = 0.8$, spending rises by 100. Find the new spending in each of the first three rounds and the total after three rounds; compare it with the limit.
\end{exercise}

\begin{exercise}[$\star\star$]\label{exo:g12:keynesian-cross:10}
From the cross in this chapter, read equilibrium output before and after the rise in spending, and compute the multiplier from the two readings.
\end{exercise}

\begin{exercise}[$\star\star$]\label{exo:g12:keynesian-cross:11}
Government spending rises by 100, and the central bank raises its rate so that investment falls by exactly 100. What is the multiplier?
\end{exercise}

\begin{exercise}[$\star\star\star$]\label{exo:g12:keynesian-cross:12}
What evidence contradicts a large multiplier, the same in all conditions? Use the opening chart and the debate.
\end{exercise}

\begin{exercise}[$\star\star\star$]\label{exo:g12:keynesian-cross:13}
In a closed economy with $c = 0.8$ and taxes that do not depend on income, government spending and taxes both rise by 100. Find the change in output.
\end{exercise}

\begin{exercise}[$\star\star\star$]\label{exo:g12:keynesian-cross:14}
An economy is at \omterm{def:g12:aggregate-demand-and-supply:potential-output}{potential output} with no idle capacity. Government spending rises. Using the \omterm{def:g12:aggregate-demand-and-supply:ad}{aggregate demand} and supply model, what happens to output and prices in the
long run, and what is the multiplier?
\end{exercise}

\begin{exercise}[$\star\star\star$]\label{exo:g12:keynesian-cross:15}
Give two reasons why open economies, and economies with a floating exchange rate, have smaller multipliers, and use the debate's estimates.
\end{exercise}

\section{Problem: Brenland's Stimulus}

\begin{problem}[{Weekend problem --- a toy multiplier and what is left of it}]
\label{pb:g12:keynesian-cross:1}
Brenland's accounts for 2001--2025 give, from lines fitted to the yearly data: a \omterm{def:g12:keynesian-cross:mpc}{marginal propensity to consume} out of disposable income $c = 0.95$; government revenue
of $t = 0.40$ of GDP in every year; a \omterm{def:g12:keynesian-cross:mpm}{marginal propensity to import} $m = 0.34$. In 2025, GDP was 393\,264 million dalens, consumption 225\,014 million and revenue 157\,306
million. The government raises its spending by 1\,\% of GDP.

\medskip
\noindent\textbf{Part I --- The toy multiplier.}
\begin{enumerate}
\item Find disposable income in 2025 and the ratio of consumption to it.
\item Check that revenue was 40\,\% of GDP in 2025.
\item Find the multiplier $1/(1 - c)$, with consumption as the only leak.
\item Find the rise in spending in millions of dalens, and the output gain with that multiplier, in millions and in per cent of GDP.
\item Find the multiplier with taxes but no imports.
\end{enumerate}

\medskip
\noindent\textbf{Part II --- Taxes and imports.}
\begin{enumerate}[resume]
\item Find the multiplier with taxes and imports.
\item Find the output gain.
\item By how much do tax receipts rise? By how much does the deficit rise?
\item By how much do imports rise?
\item Why is a gain of 20\,\% of GDP not believable even in the toy model?
\end{enumerate}

\medskip
\noindent\textbf{Part III --- The central bank and capacity.}
\begin{enumerate}[resume]
\item The central bank raises its rate so that investment and other spending fall by half the rise in government spending. Find the output gain.
\item Find the multiplier left.
\item If Brenland were at full capacity, what would the multiplier be in the long run, and what would rise instead?
\item Compare the multipliers of Parts II and III with the opening chart.
\item Which assumption of the toy model fails if households expect the extra spending to be paid by higher taxes later?
\end{enumerate}

\medskip
\noindent\textbf{Part IV --- Interpretation.}
\begin{enumerate}[resume]
\item Find the multiplier with taxes and imports if $m$ were half as large.
\item Why might a country with a fixed exchange rate see a larger multiplier than Brenland?
\item The value $c = 0.95$ is fitted on years in which income grew steadily. Why might households spend less of a one-off rise in income?
\item Which of the multipliers computed here would you report as the model's estimate, and with what warning?
\item State the result of the problem: the toy-model multiplier, and the multiplier left once imports, taxes and the monetary offset are counted.
\end{enumerate}
\end{problem}
